\section{Conclusion}
\label{sec:conclusion}
We compared openly available watermarks and fingerprints for images, audio and video as soft bindings for content provenance, under one protocol, on licence-audited public media and at calibrated false-match rates. For images, PixelSeal and a DINOv2 embedding are the strongest shippable watermark and fingerprint, and they fail on different copies; fingerprints cover more attacked audio copies, while expected-key watermark verification covers more video copies. The strongest fingerprints cannot be shipped as released, and the one open fingerprint on the C2PA list, the ISCC code, is precise but loses half of the transformed copies. Five findings concern the design of a soft-binding service more than the choice of a model. False matches are a property of the decision rule, so decoded keys must be checked against the expected payload or the registry. Catalogue variants dominate the fingerprint false-match budget, and no detection improvement from geometric verification was observed under a matched calibration false-binding constraint in the evaluated image pipelines. Neither family resists an adversary who knows the method, and unkeyed fingerprints leak content when published. Watermarking moves fingerprints, so the published asset must be the one fingerprinted. And colour conversion that differs across platforms changes decoded watermark bits.
